\chapter{Central Limit Theorem}
\section{convergence in distribution}
\begin{definition}{convergence in distribution}
我们称一个 seq of random variables $\{X_n\}$ 
converge in distribution to a random variable $X$,
写作 $X_n \overset{d}{\longrightarrow} X$, 
如果 $X$ 的 
分布函数 $F_X$ 下所有
右连续的 $x$ (即 $F_X(x) = F_X(x-)$ ), 都有 $$ 
\lim_{n \to \infty} F_{X_n}(x) = F_X(x)
 $$
\end{definition}
\begin{remark}
Convergence in distribution 的意思就是:
随着 $n$ 越来越大, $X_n$ 的分布是否越来越接近 $X$ 的分布.  
这是一个非常弱的收敛概念, 
因为它甚至\textbf{不要求 $X_n$ 和 $X$ 在同一个概率空间上定义,} 
也不要求 $X_n$ 和 $X$ 之间有任何关系. 只需要
分布函数收敛即可.

并且 notice: 并不要求对所有点 $x$ 都有 $F_{X_n}(x) \to F_X(x)$,
而是只要求对 $X$ 的分布函数 $F_X$ 下所有右连续的点 
$x$. 
这是因为如果不加这个限制, 很多直观上应该收敛的情况在数学上就会失败.

比如: 令$X_n = 1/n$, 一个 constant random variable, 
那么 $X_n$ 的分布函数 $F_{X_n}$ 是一个 step function. 
但是在 $x=0$ 处, $F_{X_n}(0) = 0$ for all $n$, 
因而 $\lim_{n \to \infty} F_{X_n}(0) = 0$. 

而 $F_X(0) = 1$. 

所以对于这种跳跃不连续点, 如果不跳过, 那么就会得出 $X_n$ 不收敛于 $X$ 的错误结论.
但是我们知道, 分布函数作为一个单调递增的函数,
是 a.e. continuous 的, 因而这些点是可以忽略掉的. 
\end{remark}


\section{Characterization of a distribution}
\subsection{moment generating function}
\begin{definition}{moment generating function}
    对于一个随机变量 $X$, 其 moment generating function (MGF) 定义为
    $$
    M_X(t) = \mathbb{E}[e^{tX}], \quad t \in \mathbb{R}
    $$
\end{definition}
\begin{remark}
为什么这个东西叫做 moment generating function 呢?
因为如果 $M_X(t)$ 在 $t=0$ 的某个 neighborhood 内存在, 
那么我们可以通过对 $M_X(t)$ 求导来得到 $X$ 的各阶矩:
\end{remark}


\begin{proposition}{}
如果 $M_X(t)$ 在 $t=0$ 的某个 neighborhood 内存在, 
则 $X$ 的 $n$ 阶矩可以表示为
    $$
    \mathbb{E}[X^n] = M_X^{(n)}(0)
    $$
其中 $M_X^{(n)}(0)$ 表示 $M_X(t)$ 在 $t=0$ 处的 $n$ 阶导数.
\end{proposition}


\subsection{characteristic function}




\section{Central Limit Theorem}
\begin{theorem}{Lindeberg-Levy Central Limit Theorem}
对于任意一个 seq of i.i.d. random variables $\{X_i\}$
 with mean $\mu$ and variance $\sigma^2 < \infty$, 
 set $S_n = X_1 + X_2 + \cdots + X_n$ for each $n$. 

我们有:
$$ 
\frac{S_n - n\mu}{\sqrt{n\sigma^2}} \overset{d}{\longrightarrow} N(0,1)
$$
\end{theorem}
\begin{remark}
LLN 告诉我们, 当 $n$ 越来越大时, $S_n/n$ 越来越接近 $\mu$. 
每个 $X_i$ 随机取样, 不论结果如何, 其平均值大概率都是 $\mu$.\\

CLT 则是进一步告诉我们这个接近的过程是如何发生的:
随着 $n$ 越来越大, $S_n$ 的分布会越来越接近一个正态分布,
且波动的幅度(标准差)数量级为 $\sqrt{n}$. 
\textbf{因而平均值 $S_n / n$ 的分布会越来越接近一个均值为 $\mu$,
 标准差为 $\sigma / \sqrt{n}$ 的正态分布}
 $$
 \frac{S_n}{n} = \mu + \frac{\sigma}{\sqrt{n}}Z, \quad Z \sim N(0,1)
$$ 
\textbf{ 随着$n$ 越来越大, 这个标准差会越来越小, 因而
其极限坍缩为一个 constant $\mu$.}\\
\end{remark}
\begin{remark}
CLT 的一个主要的应用价值, 就是可以用来估计
由多个 i.i.d. 随机因素共同作用的现象的分布. 

统计学中常用于对 confidence interval 的估计, 
以及 hypothesis testing,
通过计算如 $\mathbb{P}\left(\frac{S_n-n \mu}{\sqrt{n \sigma^2}} \in A\right)$
这种形式的概率来进行推断. \\
\end{remark}

在进行证明前, 我们先看一些 applications of CLT.
可能会帮助我们更好地理解 CLT 的意义.


\subsection{applications of CLT}
\begin{example}{(判断 coin 是否 fair)}

我们有两个 coins, 想要判断它们是否是 fair coin. 
我们可以 toss 这个 coin $n$ 次, 记录下每次 toss 的结果, 
记为 $X_1, X_2, \cdots, X_n$,
 其中 $X_i = 1$ if the $i$-th toss is heads.

现在: 观测到第一个 coin 100 次 toss 中有 38 次是 heads,
第二个 coin 100 次 toss 中有 43 次是 heads.

\begin{solution}
假设这两个 coins 是 fair coin.
那么 $X_i$ 是 i.i.d. Bernoulli random variables
 with parameter $p=0.5$, 因而
$\mu = p = \frac{1}{2}$, 
$\sigma^2 =  p(1-p) = \frac{1}{4}$. 
从而根据 CLT, 
$$ 
\mathbb{P}(S_100 < 38) 
=   
\mathbb{P}\left(\frac{S_{100} - 100 \cdot \frac{1}{2}}{\sqrt{100 \cdot \frac{1}{4}}} < \frac{38 - 100 \cdot \frac{1}{2}}{\sqrt{100 \cdot \frac{1}{4}}}\right)
\approx
\mathbb{P}\left(Z < \frac{38 - 50}{5}\right) \approx 
 \mathbb{P}(Z < -2.4) \approx 0.0082 < 0.01
$$
因而这个第一个 coin 很可能不是 fair coin.
同样的方法计算出 $\mathbb{P}\left(S_{100} \leq 43\right) \approx 0.0887$, 
因而第二个 coin 虽然也可疑不是 fair coin,
但是不如第一个 coin 可疑. 如果以 0.05 作为显著性水平, 
那么我们可以拒绝第一个 coin 是 fair coin 的假设.\\
\end{solution}
\end{example}



\begin{example}{(样本量需求的计算)}
工厂生产了一批电线, 
我们想知道它们的平均断裂强度 $\mu$ 是多少. 
遂抽取 $n$ 根电线进行测量, 
得到 $X_{1}, ..., X_{n}$, 
然后计算它们的样本平均值 $\overline{X}_{n}$ 
来估计 $\mu$. 

已知量: 强度的方差 $\sigma^{2} = 1/10$.; 
我们想估计的是 $\mu$ 的值. 
并且我们希望我们的估计是准的, 
in the sense that: 误差 $|\overline{X}_{n} - \mu|$ 不超过 0.01 的概率至少为 0.95. 

\begin{solution}
我们要达到: $\mathbb{P}(|\overline{X}_{n} - \mu| \le 1/100)$.

我们要把它转成标准的 normal distribution 的形式,
即 $Z = \frac{S_{n} - n\mu}{\sqrt{n\sigma^{2}}}$ 的形式.

于是我们把 $\overline{X}_{n} - \mu$ 写为 $\frac{S_{n}}{n} - \mu$,
然后两边同时乘 $\frac{\sqrt{n}}{\sigma}$.
右边的常数项也做相同的变换: $\frac{1/100}{\sigma / \sqrt{n}} = \frac{\sqrt{n}}{100\sigma}$.

于是原式变为:
$$
\mathbb{P}(|Z| \le \frac{\sqrt{n}}{100\sigma}) \approx 0.95
$$
对于正态分布我们知道
$$\mathbb{P}(|Z| \le x) = 2\Phi(x) - 1$$
于是想要: 
$$2\Phi\left(\frac{\sqrt{n}}{100\sigma}\right) - 1 \ge 0.95 \implies \Phi\left(\frac{\sqrt{n}}{100\sigma}\right) \ge 0.975$$
通过查表我们知道当 $\Phi(z) = 0.975$ 时,
 $z \approx 1.96$. 因此要求
 $\frac{\sqrt{n}}{100\sigma} \ge 1.96$, 
 解得至少需要 $n \ge (61.98)^{2} \approx 384.16$, 
 floor 一下得到 $n \ge 385$.
\end{solution}
\begin{remark}
这里的应用正是我们开头说的:
LLN 告诉我们, 当 $n$ 越来越大时,
 $S_n/n$ 越来越接近 $\mu$. 
每个 $X_i$ 随机取样, 不论结果如何, 
其平均值大概率都是 $\mu$; 而
CLT 则是进一步告诉我们这个接近的过程是如何发生的:
$$
 \frac{S_n}{n} = \mu + \frac{\sigma}{\sqrt{n}}Z, \quad Z \sim N(0,1)
$$ 
 随着$n$ 越来越大, 这个标准差会越来越小, 因而
其极限坍缩为一个 constant $\mu$.
\textbf{因而, 我们可以通过 CLT 来得到: 我们
要用样本均值来估计分布的真实均值的话, 
在一定的置信水平下, 需要多少样本量 $n$ 才能保证我们的估计是准的. }

\end{remark}
\end{example}



\subsection{Berry-Esseen Theorem: CLT 的收敛速度}
之前的例子中, 我们都
使用了 $\approx$ 表示:
我们直接把此时的分布近似当作了一个
 normal distribution, 来计算一些概率.

但是: 这两个分布的近似行为的本身有多么精准?

下面有一个 theorem 刻画了这件事.
\begin{theorem}{Berry-Esseen Theorem}
给定一个 seq of i.i.d. random variables $\{X_i\}$
 with mean $\mu$ and variance $\sigma^2 < \infty$, 
 以及 $\mathbb{E}[|X_i - \mu|^3]= \rho < \infty$,
 set $S_n = X_1 + X_2 + \cdots + X_n$, 
 我们有: 对于任意 $x \in \mathbb{R}$,
 $$ 
 \left|\mathbb{P}\left(\frac{S_n-n \mu}{\sqrt{n \sigma^2}} \leq x\right)-\Phi(x)\right| \leq \frac{3 \rho}{\sigma^2 \sqrt{n}}
$$
\end{theorem}
\begin{remark}
给定一个 number of samples $n$, 
我们可以给出此时的
真实概率与正态近似值之间的差值的 
一个 upper bound. \\

并且, 这个 bound 不是概率意义上的 bound,
而\textbf{是一个 deterministic bound, 
没有任何的随机性}; 并且是 
\textbf{uniform across all $x$ 的 bound.}\\

这是因为这里本身就是两个概率的值之间的差,
而不是样本之间的差. 因而其实这是一个非常强大的
结论, 表面\textbf{我们用正态分布近似估计出的概率
和真实的概率之间的差值是 determinstically 
controllable 意义上非常小的.}\\

我们马上可以想到: 对于
控制样本数量来达到某个置信水平的问题, 
这是一个非常有用的工具, 因为它告诉我们, 当 $n$ 足够大时,
我们用正态分布近似估计出的概率和真实的概率之间的
差值是非常小并且可以控制的.\\
并且, 这个控制应该是比较精准的. 
因为相比其他的控制方法:
\textbf{Chebyshev 只需要 first, second order 
moments 的信息; Hoeffding 
只需要有界性的假设, 甚至不需要知道
分布的任何 moment 信息; 而 
Berry-Esseen 需要 third order moment 的信息, 
因此它应该能够给一个更精准的控制.}

\end{remark}
\begin{proof}
//TODO:
\end{proof}

\begin{example}
一个工厂生产电子元件, 
每个原件有概率是有缺陷的. 

令 
$$
X_i= \begin{cases}1, & \text { if the } i \text { th tested component is defective }, \\ 0, & \text { otherwise } .\end{cases}
$$
并假设 i.i.d. with parameter
$$ 
\mathbb{P}(X_i = 1) = p
 $$
 其中 $p$ 是一个未知的参数.

工厂方面表示这个 process 是 under control 的,
并给出了假设: 
$H_0: p \leq 0.02$.

为了验证这个假设, 
一个 quality control manager 
需要从生产线上需要随机 $n$ 个元件进行测试, 
并估计出 $p$ 的值 by 样本均值:
$$
\hat{p}_n:=\frac{1}{n} \sum_{i=1}^n X_i
$$
我们希望这个估计的误差最多为 0.005,
并且这个误差的置信度至少为 0.99, 
即我们希望:
$$ 
\mathbb{P}\left(\left|\hat{p}_n-p\right|>0.005\right) \leq 0.01 
$$

那么我们至少需要多少样本量 $n$ 来达到这个要求呢?


\begin{solution}
首先计算样本均值 $\hat{p}_n$ 的 mean 和 variance: 
$$ 
\mathbb{E}\left[\hat{p}_n\right]=p, \quad \operatorname{Var}\left(\hat{p}_n\right)=\frac{1}{n^2} \sum_{i=1}^n \operatorname{Var}\left(X_i\right)=\frac{p(1-p)}{n} 
$$
此时我们有三个办法:
\begin{itemize}
\item 
办法1: \textbf{Chebyshev's inequality}.
$$ 
\mathbb{P}\left(\left|\hat{p}_n-p\right| \geq \varepsilon\right) \leq \frac{\operatorname{Var}\left(\hat{p}_n\right)}{\varepsilon^2}=\frac{p(1-p)}{n \varepsilon^2} 
$$
由于 $p(1-p) \leq 1/4$ for any $p \in [0,1]$,
式子可以进一步控制为
$$ 
\mathbb{P}\left(\left|\hat{p}_n-p\right| \geq \varepsilon\right) \leq \frac{1}{4 n \varepsilon^2} 
$$
我们希望此概率不超过 0.01,
并且\ $\varepsilon = 0.005$, 那么进一步得到需要
$$ 
\frac{1}{4 n(0.005)^2} \leq 0.01 
$$
解得至少需要 $n\geq 1000000$.

\item 
办法2: \textbf{Hoeffding's inequality}.\\
由于 $0\leq X_i \leq 1$, 
我们可以直接使用 Hoeffding's inequality 来控制: 
$$ \mathbb{P}\left(\left|\hat{p}_n-p\right| \geq \varepsilon\right) \leq 2 e^{-2 n \varepsilon^2}
 $$
我们 require 了 $2 e^{-2 n(0.005)^2} \leq 0.01$, 
因而可以解得
$$ n \geq \frac{|\ln (0.005)|}{2(0.005)^2} \approx 106,000 $$

\item 
办法3: \textbf{Berry-Esseen Theorem}.\\
首先我们假设报告的 $p$ 的值 0.02 是正确的,
从而可以计算出:
$$
\mathbb{E}\left[X_1\right]=p, \quad \sigma^2:=\operatorname{Var}\left(X_1\right)=p(1-p)=0.0196, \quad \sigma=0.14
$$
并且, 由于 $\left|X_1-p\right| \leq 1$, 
我们知道偏度 $\rho:=\mathbb{E}\left[\left|X_1-p\right|^3\right]<\infty$.

因而可以应用 Berry-Esseen Theorem. 
我们希望:
$$ 
\mathbb{P}\left(\left|\hat{p}_n-p\right| \leq \varepsilon\right)=\mathbb{P}\left(\left|S_n-n p\right| \leq n \varepsilon\right) \geq 0.99, \quad \varepsilon=0.005 
$$

等价于
$$ \mathbb{P}\left(\left|\frac{S_n-n p}{\sigma \sqrt{n}}\right| \leq \frac{\varepsilon \sqrt{n}}{\sigma}\right) \geq 0.99
 $$
Berry-Esseen Theorem 可以得到:
对于任意的 $x > 0$,
有 
$$
\mathbb{P}\left(\left|\frac{S_n-n p}{\sigma \sqrt{n}}\right| \leq x\right) \geq 2 \Phi(x)-1-\frac{6 \rho}{\sigma^2 \sqrt{n}}
$$
因而我们需要选择 $n$ s.t. 
$$ 2 \Phi\left(\frac{\varepsilon \sqrt{n}}{\sigma}\right)-1-\frac{6 \rho}{\sigma^2 \sqrt{n}} \geq 0.99 
$$
忽略掉 $\frac{6 \rho}{\sigma^2 \sqrt{n}}$ 这个很小的项
(计算可以再精细地选择 $n$), 
我们也可以得到:

$$ \frac{\varepsilon \sqrt{n}}{\sigma} \geq z_{0.995} \approx 2.576 $$
因而至少需要 $n \geq 5200$. 
\end{itemize}
我们可以看到
\begin{itemize}
    \item 通过 Chebyshev's inequality 来控制, 我们需要的样本量是 100 万级别的; 
    \item 通过 Hoeffding's inequality 来控制, 我们需要的样本量是 10 万级别的; 
    \item 通过 Berry-Esseen Theorem 来控制, 我们需要的样本量是 5000 级别的, 显然更为高效. 
\end{itemize}
不过, 这个结果是基于我们假设 $p=0.02$ 
的前提下得到的.
Chebyshev 的逻辑 (保守估计): 我不相信厂家的任何话, 我要建立一个无论真实 $p$ 是多少都绝对成立的置信区间. 所以我必须用最差的情况 $p=0.5$ 来算 $n$.
CLT/Berry-Esseen 的逻辑 (假设检验):  
厂家宣称 $p \le 0.02$. 在统计学中, 
我们通常先假设这个宣称是正确的 (即 $H_0$ 假设).
如果我们用在这个假设下得出的需要的样本量
 $n=5200$ 去测, 发现结果远超 $0.02$, 那我们就直接推翻厂家的说法. 

\end{solution}
\end{example}




\subsection{proof of CLT}
